% Unit 30 — Gaussian and Dirichlet
\input{technique-note.tex}
\title{Integration Bee --- Unit 30:\\ Gaussian and Dirichlet Integrals}
\author{UK Uni Integration Bee prep\\ Named integrals, with the reductions that stay elementary}
\date{}
\begin{document}
\maketitle
\section*{Status}
Tags: \texttt{yang:5.3}, \texttt{yang:5.4}; techniques: \texttt{advanced-gaussian}, \texttt{dirichlet}.
Two named values do the heavy lifting: $\int_{-\infty}^{\infty}e^{-x^2}\,dx=\sqrt{\pi}$ and $\int_0^{\infty}(\sin x)/x\,dx=\pi/2$. Everything else in the drill set is a substitution, a parameter derivative, or one parts step on top of those values. Yang \S\S5.3--5.4 records the contest forms. The evaluations below are the usual public arguments.
\section{Official anchors}
\begin{itemize}
\item Unit 11 for the first Gaussian cousins. Unit 34 for a parameter derivative.
\item Syllabus \S2 special functions, as names. Vector calculus is not required for Round 1; the polar sketch below is optional background.
\item Yang Topics Guide 2nd Ed \S\S5.3 Advanced Gaussian, 5.4 Dirichlet Integral.
\end{itemize}
\section{Learning goals}
\begin{enumerate}
  \item Quote $\int_0^{\infty}e^{-x^2}\,dx=\sqrt{\pi}/2$, then scale. Do not invent an elementary antiderivative of $e^{-x^2}$.
  \item Integrate $x e^{-x^2}$ by substituting $u=x^2$. That one is elementary.
  \item Quote $\int_0^{\infty}(\sin x)/x\,dx=\pi/2$, then change variables or integrate by parts.
\end{enumerate}

\section{The Gaussian}
Evenness gives
\[
\int_0^{\infty}e^{-x^2}\,dx=\frac12\int_{-\infty}^{\infty}e^{-x^2}\,dx=\frac{\sqrt{\pi}}{2},
\]
which is \texttt{U30-D01}. The full line is \texttt{U30-D03}. One evaluation of the square, optional here, is the polar computation
\[
\Bigl(\int_{-\infty}^{\infty}e^{-x^2}\,dx\Bigr)^2=\int_0^{2\pi}\int_0^{\infty}e^{-r^2}r\,dr\,d\theta=\pi.
\]
For $a>0$, the substitution $u=x\sqrt{a}$ yields $\int_0^{\infty}e^{-a x^2}\,dx=\frac12\sqrt{\pi/a}$. \texttt{U35-D01} is $a=2$.

\texttt{U30-D02} is elementary:
\[
\int_0^{\infty}x e^{-x^2}\,dx=\frac12.
\]
\texttt{U30-D04} differentiates the scaled Gaussian in $a$ and sets $a=1$:
\[
\int_0^{\infty}x^2 e^{-x^2}\,dx=\frac{\sqrt{\pi}}{4}.
\]
The differentiation is the Leibniz rule from unit 34, taken as a limit of compact intervals.

\section{The Dirichlet integral --- \texttt{U30-D05}}
For $a>0$ set $I(a)=\int_0^{\infty}e^{-ax}(\sin x)/x\,dx$. Differentiating in $a$ under the integral gives
\[
I'(a)=-\int_0^{\infty}e^{-ax}\sin x\,dx=-\frac{1}{a^2+1}.
\]
$I(a)\to 0$ as $a\to\infty$, so $I(a)=\pi/2-\arctan a$. The limit $a\to 0^+$ is $\pi/2$. That is the stored answer of \texttt{U30-D05}. The switch of derivative and integral is a limit; unit 34 states the compact case.

\section{One parts step --- \texttt{U30-D06}}
\[
\int_0^{\infty}\frac{\sin^2 x}{x^2}\,dx.
\]
Parts with $u=\sin^2 x$ and $dv=x^{-2}\,dx$ kills the boundary term in the limit and leaves $\int_0^{\infty}\sin(2x)/x\,dx$. The substitution $u=2x$ turns that into $\int_0^{\infty}(\sin u)/u\,du=\pi/2$.

\texttt{U30-D07} and \texttt{U30-D08} are $\int_0^{\infty}e^{-x}\sin x\,dx$ and the cosine partner. Both equal $1/2$. They are one-line parts, or the real and imaginary parts of $\int_0^{\infty}e^{-x}e^{ix}\,dx$ in unit 32. They do not need the Dirichlet value.

\section{Drill pointers}
See \texttt{drills.tex} / \texttt{drills.md}. Feed: \texttt{verify/drills\_unit30.json}.

\section*{On a miss}
Teach this file. Say which named value you used, then the scaling or the parts step. Do not invent an answer that is not in \texttt{drills.md}.
\end{document}
